% Original applied explanation; see Chapter04/SOURCES.md.
\subsubsection{একই বিজ্ঞপ্তিতে div, span ও CSS-এর ভূমিকা}
কলেজের একটি বিজ্ঞপ্তির শিরোনাম, তারিখ ও লেখা একত্রে ধরতে \texttt{div} ব্যবহার করা যায়। অনুচ্ছেদের শুধু \textbf{শেষ তারিখ} অংশে আলাদা রীতি প্রয়োগ করতে \texttt{span} উপযোগী। CSS-এর \texttt{color} লেখার রং এবং \texttt{background-color} পটভূমির রং নির্ধারণ করে।

\texttt{div} ও \texttt{span} নিজেরা বিজ্ঞপ্তি বা মেনুর নির্দিষ্ট অর্থ জানায় না। প্রধান বিষয়বস্তুর জন্য \texttt{main}, নেভিগেশনের জন্য \texttt{nav}-এর মতো অর্থবোধক অংশ বেছে নেওয়া যায়। কোন ট্যাগ কী অর্থ দেয় এবং CSS সেটিকে কীভাবে দেখায়—দুটি প্রশ্ন আলাদা করে ভাববে।
